\subsection{\texorpdfstring{有界复测度 \(^{6}\)}{有界复测度 (note 6)}}

对每个复测度 \(m\) （ \(\mathrm{T}\) 上的复测度），令

\[
\| m \| = \sup _ {\| f \| \leqslant 1,   f \in \mathcal {K} _ {\mathbf {C}} (\mathrm{T})} | m (f) |.
\]

若 \(\|m\| < +\infty\) ，则称 m 有界；这等价于说 m 在赋以一致收敛拓扑的 \(\mathcal{K}_{\mathbf{C}}(\mathrm{T})\) 上连续，从而可以延拓为一个连续线性形式（范数为 \(\|m\|\) ），它定义在无穷远处趋于 0 的连续复函数所成的 Banach 空间 \(\overline{\mathcal{K}_{\mathbf{C}}(\mathrm{T})}\) 上。

引理 5. --- 设 m 是 T 上的复测度，f 是 m-可积的复函数。则 \(\int |f| d|m| = \sup \left| \int f h dm \right|\) ，其中 h 跑遍 \(\mathcal{K}_{\mathbf{C}}(\mathbf{T})\) 中使 \(|h(t)| \leqslant 1\) （对一切 \(t \in T\) ）成立的函数所成之集。

若 \(m = g \cdot |m|\) ，则 \(\int |f| d|m| = \int |fg| d|m|\) 且 \(\int f h dm = \int f gh d|m|\) 。令 \(\zeta(t) = 0\) （当 \(f(t)g(t) = 0\) 时），令 \(\zeta(t) = \frac{f(t)g(t)}{|f(t)g(t)|}\) （当 \(f(t)g(t) \neq 0\) 时）； \(\zeta\) 是 \(|m|\) -可测的，于是对每个 \(\varepsilon > 0\) ，存在 T 的紧子集 K，使 \(\int_{T-K} |f| d|m| \leqslant \varepsilon\) ， \(\zeta\) 在 K 上的限制连续，且在 K 上 \(|\zeta(t)| = 1\) 。因此，借助 Urysohn 定理，存在一个定义在 T 上、取复值的连续函数 \(\zeta_1\) ，使得在 K 上 \(\zeta_1 = \zeta\) ，且 \(|\zeta_1(t)| \leqslant 2\) 及 \(\zeta_1(t) \neq 0\) （对每个 \(t \in T\) ）；令 \(h(t) = \zeta_1(t)/|\zeta_1(t)|\) ，便知 h 在 T 上连续，在 K 上与 \(\zeta\) 一致，且满足 \(|h(t)| = 1\) （对所有 \(t \in T\) ）。最后，令 u 为 T 到 {[}0,1{]} 中的连续映射，它在 K 上等于 1 且具有紧支集；令 \(h_1 = h^{-1}u\) ，我们有

\[
\left| \int f h _ {1} d m - \int | f | d | m | \right| \leqslant 2 \int_ {\mathrm{T-K}} | f | d | m | \leqslant 2 \varepsilon ,
\]

这就证明了引理。

命题 13. --- 设 m 是复测度， \(\mu = |m|\) 。为使 m 有界，必须且只需 \(\mu\) 有界，且这时 \(\|m\| = \|\mu\|\) 。

我们有 \(m = g \cdot \mu\) ，其中 \(g\) 是 \(\mu\) -可测的，且 \(|g(t)| = 1\) （对所有 \(t \in \mathbf{T}\) ）。若 \(\mu\) 有界，则对每个函数 \(f \in \mathcal{H}_{\mathbf{C}}(\mathbf{T})\) ，

\[
| m (f) | = \left| \int f g d \mu \right| \leqslant \mathrm{N} _ {\infty} (f g) \| \mu \| = \| f \| \cdot \| \mu \|,
\]

因此 m 有界且 \(\|m\|\leqslant\|\mu\|\) 。若 m 有界，则对每个 \(f\in\mathcal{K}_{\mathbf{C}}(\mathrm{T})\) ，考虑到引理 5，我们有

\[
| \mu (f) | \leqslant \| f \| \cdot \| m \|,
\]

因此 \(\mu\) 有界且 \(\| \mu \| \leqslant \| m\|\) 。命题得证。

推论. --- 设 m 是有界复测度。则每个函数 \(\mathbf{f} \in \mathcal{L}_{\mathrm{F}}^{\infty}(m)\) 都是 m-可积的，且 \(\left|\int f dm\right| \leqslant N_{\infty}(\mathbf{f}) \|m\|\) 。

因为， \(\mathbf{f}\) 是 \(m\) -可测的，且令 \(\mu = |m|\) ，我们有

\[
\int^ {*} | \mathbf {f} | d \mu \leqslant \mathrm{N} _ {\infty} (\mathbf {f}) \| \mu \| = \mathrm{N} _ {\infty} (\mathbf {f}) \| m \|,
\]

因此 \(\mathbf{f}\) 是 \(|m|\) -可积的（Ch. IV, §5, No.~6, 定理 5），且

\[
\left| \int \mathbf {f} d m \right| \leqslant \int | \mathbf {f} | d \mu \leqslant \mathrm{N} _ {\infty} (\mathbf {f}) \| m \|.
\]

\begin{enumerate}
\def\labelenumi{\arabic{enumi}.}
\setcounter{enumi}{9}
\tightlist
\item
  复测度的像；诱导复测度；复测度的乘积 \(^{7}\)
\end{enumerate}

设 \(m\) 是 \(T\) 上的复测度， \(\pi\) 是 \(T\) 到局部紧空间 \(X\) 中的映射。我们说 \(\pi\) 是 \(m\) -逆紧的，意思就是 \(\pi\) 是 \(|m|\) -逆紧的（Ch. V, §6, No.~1, 定义 1）；于是立即得出，对每个函数 \(f \in \mathcal{H}_{\mathbf{C}}(X)\) ，函数 \(f \circ \pi\) 本质 \(m\) -可积，且

\[
\left| \int (f \circ \pi) d m \right| \leqslant \int | f \circ \pi | d | m | = \int | f | d (\pi (| m |)),\tag{8}
\]

因此映射 \(f \mapsto \int (f \circ \pi) dm\) 在 \(\mathcal{K}_{\mathbf{C}}(\mathbf{X})\) 上连续，换言之，它是 \(\mathbf{X}\) 上的一个复测度，记作 \(\pi(m)\) ，称为 \(m\) 在 \(\pi\) 下的像。此外，由 (8) 得出 \(|\pi(m)| \leqslant \pi(|m|)\) 。若 \(m\) 与 \(m'\) 是 \(T\) 上的两个复测度，且 \(\pi\) 既 \(m\) -逆紧又 \(m'\) -逆紧，则 \(\pi\) 是 \((\lambda m + \lambda'm')\) -逆紧的（对任意复标量 \(\lambda, \lambda'\) ），且 \(\pi(\lambda m + \lambda'm') = \lambda\pi(m) + \lambda'\pi(m')\) 。

设 Y 是 T 的局部紧子空间。对每个函数 \(f \in \mathcal{K}_{\mathbf{C}}(\mathbf{Y})\) ，T 上的函数 \(f'\) （由 \(f'(t) = f(t)\) （当 \(t \in Y\) 时）与 \(f'(t) = 0\) （当 \(t \notin Y\) 时）定义）是 m-可积的（Ch. IV, §5, No.~7）；立即可以看出，映射 \(f \mapsto \int f' dm\) 是 Y 上的一个复测度，称为 m 在 Y 上诱导的测度，记作 \(m_{Y}\) 。若 \(m = g \cdot |m|\) ，则显然 \(m_{Y} = g_{Y} \cdot |m|_{Y}\) ，其中 \(g_{Y}\) 是函数 g 在 Y 上的限制，它关于 \(|m|_{Y}\) 局部可积（Ch. V, §7, No.~1）；此外，由于 \(|g_{Y}| = 1\) 关于 \(|m|_{Y}\) 局部几乎处处成立（Ch. V, §7, No.~1, 命题 1 的推论 1），我们有 \(|m_{Y}| = |m|_{Y}\) 。

设 T 与 T' 是两个局部紧空间，m（相应地， \(m'\) ）是 T（相应地，T'）上的复测度。写 \(m = g \cdot |m|\) 与 \(m' = g' \cdot |m'|\) 。函数 \(g \otimes g'\) 关于正测度 \(|m| \otimes |m'|\) 在 T × T' 上局部可积（Ch. V, §8, No.~5, 命题 10），且立即可以验证：若 g（相应地， \(g'\) ）被换成一个函数 \(g_1\) （相应地， \(g'_1\) ），它与 g（相应地， \(g'\) ）关于 \(|m|\) （相应地， \(|m'|\) ）局部几乎处处相等，则 \(g_1 \otimes g'_1\) 等于 \(g \otimes g'\) （关于 \(|m| \otimes |m'|\) 局部几乎处处）。因此，T × T' 上的复测度 \((g \otimes g') \cdot (|m| \otimes |m'|)\) 只依赖于 m 与 \(m'\) ；把它记作 \(m \otimes m'\) ，称为 m 与 \(m'\) 的乘积测度。由于 \(|g \otimes g'| = 1\) 关于 \(|m| \otimes |m'|\) 局部几乎处处成立（Ch. V, §8, No.~2, 命题 4），我们有 \(|m \otimes m'| = |m| \otimes |m'|\) 。

读者容易验证，Ch. V 中已证明的关于正测度的像、正测度诱导的测度以及正测度的乘积的所有命题，除了其中出现上积分或本质上积分者外，当把正测度换成任意复测度时仍然成立。

最后，像 §1 中那样，对复测度 m 定义标量本质 m-可积函数的概念；为使函数 f 具有这一性质，必须且只需 f 关于 \(|\mu_{1}|\) 与 \(|\mu_{2}|\) 标量本质可积（其中 \(\mu_{1}\) 与 \(\mu_{2}\) 是 m 的实部与虚部），且这时 \(\int f dm = \int f d\mu_{1} + i \int f d\mu_{2}\) 。 §1 的结果向复测度的移植工作留给读者。

